\section{All-rank localization of ordered finite-difference sums}
\label{sec:localization}

The incoming finite-difference report proves determinant identities for
symmetrized assignments of orders and shifts, and explicitly asks which
coordinates are needed for an individual ordered assignment
\cite[Sections 3 and 9]{repo:finite}.  We give a finite construction for
arbitrary rank.  Its central fact is a support cancellation: after the
differences are transferred to the ordered Vandermonde polynomial, two
isolated free labels annihilate the coefficient.  The construction retains
the individual assignment of every order and shift.

\subsection{The ordered sum and its exact transported coefficient}

Fix an integer $r\ge2$ and put $d=2r-2$.  For
$\mathbf s=(s_1,\ldots,s_r)\in\C^r$ and
$\Rea a_i>0$, define
\begin{align}
 w_i(n)&=\frac1{d!}\sum_{j=0}^{d}(-1)^j\binom dj
                       (n+a_i+j)^{-s_i},
                       \label{loc:weights}\\
 \mathcal Z_r(\mathbf s;\mathbf a)
   &=\sum_{0\le n_1<\cdots<n_r}
          \Delta(\mathbf n)^2\prod_{i=1}^r w_i(n_i),
                       \label{loc:ordered-sum}
\end{align}
where $\Delta(\mathbf x)=\prod_{i<v}(x_v-x_i)$.
All powers use the principal logarithm in the right half-plane.  We shall
distinguish two domains:
\[
 \Omega=\{\Rea s_i>1,\ \Rea a_i>0\text{ for every }i\},
 \qquad
 \Omega_0=\{\Rea s_i>d+1,\ \Rea a_i>0\text{ for every }i\}.
\]

\begin{theorem}[Finite ordered reduction in every rank]
\label{loc:main}
The series \eqref{loc:ordered-sum} converges absolutely and locally normally
on $\Omega$.  On $\Omega_0$ it is an explicit finite sum of expressions
\begin{equation}
 \sum_{0\le N_1<\cdots<N_q}
    Q(N_1,\ldots,N_q)
    \prod_{\ell=1}^q\prod_{i\in B_\ell}
                      (N_\ell+A_i)^{-s_i},
 \qquad \Rea A_i>0,
 \label{loc:reduced-form}
\end{equation}
multiplied, where necessary, by finite products of shifted powers.
Here $Q$ is a polynomial, the label sets $B_\ell$ are consecutive blocks,
and
\[
 q\le\left\lfloor\frac{r+1}{2}\right\rfloor.
\]
The contribution anchored at the lower boundary has at most
$\lfloor r/2\rfloor$ free indices.  Equations
\eqref{loc:sector-polynomial}, \eqref{loc:free-sector}, and
\eqref{loc:anchored-sector} below give all coefficients and shifts by
finite sums.  Every expression in this construction converges absolutely
on $\Omega_0$.  Their collective sum has a holomorphic continuation to
$\Omega$, equal there to \eqref{loc:ordered-sum}.
\end{theorem}

The number $q$ counts summation variables in the displayed representation.
The theorem makes no assertion that this number is minimal, or that any
special values obtained from different sectors are arithmetically
independent.

For the convergence statement, the finite-difference integral is
\begin{equation}
 w_i(n)=\frac{(s_i)_d}{d!}
       \int_{[0,1]^d}
       (n+a_i+u_1+\cdots+u_d)^{-s_i-d}\dd\mathbf u.
 \label{loc:difference-integral}
\end{equation}
The sign is positive because $d$ is even.  On compact subsets of the
parameter domain this gives
$|w_i(n)|\le C(1+n)^{-\Rea s_i-d}$.  For nonnegative integer nodes,
\begin{equation}
 |\Delta(\mathbf n)|^2\le\prod_{i=1}^r(1+n_i)^d.
 \label{loc:vandermonde-bound}
\end{equation}
The resulting product majorant proves normal convergence on $\Omega$.
If the differences in \eqref{loc:weights} are expanded before summation,
each factor is only bounded by $C(1+n_i)^{-\Rea s_i}$; the same argument
then proves absolute convergence on $\Omega_0$.  This is the initial
domain for the coefficient transport.

On integer tuples define the chamber polynomial
\begin{equation}
 P(\mathbf x)=\ind_{\{0\le x_1<\cdots<x_r\}}
                                      \Delta(\mathbf x)^2
 \label{loc:chamber}
\end{equation}
and the finite stencil
\begin{equation}
 C_r(\mathbf k)=
 \sum_{\mathbf j\in\{0,\ldots,d\}^r}
       (-1)^{j_1+\cdots+j_r}\prod_{i=1}^r\binom d{j_i}
       P(\mathbf k-\mathbf j).
 \label{loc:transport-coefficient}
\end{equation}
Expanding \eqref{loc:weights}, setting $k_i=n_i+j_i$, and collecting the
common shifted powers gives the exact identity
\begin{equation}
 \mathcal Z_r(\mathbf s;\mathbf a)
 =\frac1{(d!)^r}\sum_{k_1,\ldots,k_r\ge0}
               C_r(\mathbf k)\prod_{i=1}^r(k_i+a_i)^{-s_i}
 \qquad\text{on }\Omega_0.
 \label{loc:transport}
\end{equation}
Expanded absolute convergence justifies every regrouping.  There is no
unaccounted upper endpoint: the regrouping is performed on the full
absolutely convergent series.  Also, $C_r(\mathbf k)=0$ if any $k_i<0$,
since every nonzero term in \eqref{loc:transport-coefficient} has
$k_i=x_i+j_i\ge0$.

\subsection{Collision blocks and the two-singleton cancellation}

The shifted tuple $\mathbf k$ need not be increasing.  If
$C_r(\mathbf k)\ne0$, however, some stencil term has a strictly increasing
tuple $\mathbf k-\mathbf j$.  Necessarily
\begin{equation}
 k_{i+1}-k_i\ge1-d\qquad(1\le i<r).
 \label{loc:gap-lower}
\end{equation}
Partition the labels $1,\ldots,r$ into maximal consecutive blocks by
placing a cut precisely when
\begin{equation}
 k_{i+1}-k_i>d.
 \label{loc:cut}
\end{equation}
The internal adjacent gaps therefore belong to the finite set
$\{1-d,\ldots,d\}$.  Across a cut, every stencil assignment remains
strictly ordered, because
$k_{i+1}-j_{i+1}-k_i+j_i\ge d+1-d=1$.

The first block is called \emph{anchored} when $k_1<d$; every other block
is free.  If $k_1\ge d$, all blocks are free.  Only the first block can
be anchored.  Indeed, the chamber condition in \eqref{loc:chamber}
requires just $x_1\ge0$ together with strict order, which already implies
nonnegativity of the later coordinates.

\begin{lemma}[Two free singleton blocks give zero]
\label{loc:singleton}
If the block decomposition of $\mathbf k$ contains two free blocks of
size one, then $C_r(\mathbf k)=0$.
\end{lemma}

\begin{proof}
Let $i$ be one free singleton, and fix every stencil variable except
$j_i$.  All chamber indicators are independent of $j_i$: its neighboring
cuts remain ordered for every shift, and it cannot change the lower
boundary condition.  As a polynomial in $x_i=k_i-j_i$, the Vandermonde
square has degree $d=2r-2$, with leading coefficient
\[
 \Delta(\mathbf x_{\widehat i})^2
       =\prod_{\substack{p<v\\p,v\ne i}}(x_v-x_p)^2.
\]
Consequently
\[
 \sum_{j_i=0}^d(-1)^{j_i}\binom d{j_i}\Delta(\mathbf x)^2
       =d!\,\Delta(\mathbf x_{\widehat i})^2.
\]
Let $v$ be the second free singleton.  Its indicators are still
independent of $j_v$, but the remaining polynomial has degree
$2(r-2)=d-2$ in $x_v$.  Its $d$th backward difference is zero.  Summing
over the other stencil variables preserves this cancellation.
\end{proof}

At most one surviving free block can therefore be a singleton.  If
there are $q$ free blocks and no anchor, they use at least $2q-1$ labels;
hence $q\le\lfloor(r+1)/2\rfloor$.  If an anchor is present, it uses at
least one additional label, giving $r\ge2q$ and
$q\le\lfloor r/2\rfloor$.  This proves the support part of
Theorem~\ref{loc:main}.  We now turn it into a finite formula without
leaving any block-coordinate choices implicit.

\subsection{An explicit finite construction of every sector}

Choose a consecutive block composition
$\mathcal B=(B_1,\ldots,B_p)$ of $\{1,\ldots,r\}$, and write
$b_\ell=|B_\ell|$ and $\ell(i)$ for the block containing $i$.  In each
block choose integer offsets $u_i$ with first offset zero and
\begin{equation}
 u_{i+1}-u_i\in\{1-d,\ldots,d\}
                    \quad(i,i+1\in B_\ell).
 \label{loc:offsets}
\end{equation}
There are finitely many choices.  Put
\[
 k_i=t_{\ell(i)}+u_i,
 \qquad c_\ell=d+u_{\max B_\ell},
\]
and introduce the finite internal indicator
\[
 \varepsilon_{\mathcal B,\mathbf u}(\mathbf j)
 =\prod_{\ell=1}^p\prod_{\substack{i,i+1\in B_\ell}}
       \ind_{\{u_{i+1}-u_i+j_i-j_{i+1}\ge1\}}.
\]
Discard a pattern if this indicator is zero for every stencil tuple.
For every retained pattern, an admissible tuple gives
\begin{equation}
 u_i\ge i-\min B_\ell-d\quad(i\in B_\ell),
 \qquad c_\ell\ge b_\ell-1\ge0.
 \label{loc:offset-feasibility}
\end{equation}
To see the first inequality, sum the internal strict inequalities from
the first label of the block to $i$, then use $j_i-j_{\min B_\ell}\ge-d$.
The second is its instance at the last label.

Define the following polynomial in the block bases:
\begin{align}
 \mathcal C_{\mathcal B,\mathbf u}(\mathbf t)
 ={}&\sum_{\mathbf j\in\{0,\ldots,d\}^r}
       (-1)^{|\mathbf j|}\prod_{i=1}^r\binom d{j_i}
       \varepsilon_{\mathcal B,\mathbf u}(\mathbf j)
       \prod_{i<v}
       \bigl(t_{\ell(v)}-t_{\ell(i)}+u_v-u_i-j_v+j_i\bigr)^2.
 \label{loc:sector-polynomial}
\end{align}
It has integer coefficients and satisfies
\begin{equation}
 \deg_{t_\ell}\mathcal C_{\mathcal B,\mathbf u}
                         \le 2b_\ell(r-b_\ell).
 \label{loc:sector-degree}
\end{equation}
Pairs of labels in the same block do not depend on its base; each of
the $b_\ell(r-b_\ell)$ cross-block pairs contributes at most degree two.

For an unanchored sector the base conditions are exactly
\[
 t_1\ge d,\qquad t_{\ell+1}>t_\ell+c_\ell.
\]
Set
\begin{equation}
 e_1=d,\qquad e_\ell=d+\sum_{v<\ell}c_v,
 \qquad t_\ell=N_\ell+e_\ell.
 \label{loc:free-bases}
\end{equation}
The conditions become $0\le N_1<\cdots<N_p$.  Moreover,
$e_\ell+u_i\ge0$ by \eqref{loc:offset-feasibility}, so every shifted
argument stays in the right half-plane.  This sector contributes exactly
\begin{equation}
 \frac1{(d!)^r}
 \sum_{0\le N_1<\cdots<N_p}
       \mathcal C_{\mathcal B,\mathbf u}(\mathbf N+\mathbf e)
       \prod_{\ell=1}^p\prod_{i\in B_\ell}
                    (N_\ell+e_\ell+u_i+a_i)^{-s_i}.
 \label{loc:free-sector}
\end{equation}
The stencil polynomial is exactly $C_r(\mathbf k)$ on this sector:
cross-block order and the lower-bound condition are automatic, leaving
precisely the indicators in \eqref{loc:sector-polynomial}.

For an anchored sector choose the finite integer
$h=k_1\in\{0,\ldots,d-1\}$ and set $t_1=h$.  Let
$\mathcal C_{\mathcal B,\mathbf u;h}(t_2,\ldots,t_p)$ be
\eqref{loc:sector-polynomial} with $t_1=h$ and the additional indicator
$\ind_{\{j_1\le h\}}$ inserted in the summand.  Discard patterns with
no admissible stencil tuple under this extra constraint.  If $p\ge2$,
set
\begin{equation}
 e_2=h+c_1+1,\qquad
 e_\ell=e_2+\sum_{v=2}^{\ell-1}c_v\quad(\ell\ge3),
 \qquad t_\ell=N_\ell+e_\ell.
 \label{loc:anchored-bases}
\end{equation}
The free variables satisfy $0\le N_2<\cdots<N_p$.  Feasibility of the
anchored block gives $h+u_i\ge0$ for its labels, and in particular
$e_2=d+h+u_{\max B_1}+1\ge d+1$.  Together with
\eqref{loc:offset-feasibility}, this also ensures positive real parts for
all the free shifted arguments.  The anchored sector is
\begin{align}
 &\frac1{(d!)^r}
   \prod_{i\in B_1}(h+u_i+a_i)^{-s_i}
   \sum_{0\le N_2<\cdots<N_p}
     \mathcal C_{\mathcal B,\mathbf u;h}
                      (N_2+e_2,\ldots,N_p+e_p)
 \notag\\[-2pt]
 &\hspace{44mm}\times
       \prod_{\ell=2}^p\prod_{i\in B_\ell}
                   (N_\ell+e_\ell+u_i+a_i)^{-s_i}.
 \label{loc:anchored-sector}
\end{align}
For $p=1$, the sum and the free product are absent, so the contribution
is a finite scalar expression.

Every possible nonzero coefficient in \eqref{loc:transport} belongs to
exactly one of these sectors: its cuts determine the composition, the
first coordinate in each block determines the base, and subtraction
determines the offsets.  Conversely, the indicated inequalities give
exactly those cuts and internal gaps.  There are finitely many
compositions, offsets, and anchored values $h$.  Summing
\eqref{loc:free-sector} and \eqref{loc:anchored-sector} over them, and
discarding the zero sectors of Lemma~\ref{loc:singleton}, is therefore
the promised finite construction.

For completeness, each separated sector converges absolutely on
$\Omega_0$.  If a free block has $b$ labels, its polynomial degree in
the associated variable is at most $2b(r-b)$, while the sum of the real
parts of its denominator exponents is greater than $b(d+1)$.  Their
difference satisfies
\[
 b(d+1)-2b(r-b)=b(2b-1)\ge1,
\]
with strict improvement supplied by $\Rea s_i>d+1$.  A product
majorant over unrestricted nonnegative free variables is summable.
The same bounds apply locally uniformly to the parameters.  This proves
the convergence assertions and completes the proof of
Theorem~\ref{loc:main}.

If a representation without polynomial numerators is preferred, choose
one denominator $(N+A)^{-s}$ in each free block.  Expanding
$N^k=((N+A)-A)^k$ removes any monomial numerator by finitely many integer
decrements of that exponent.  The preceding degree estimate proves
absolute convergence of these terms on $\Omega_0$ as well, and no new
summation variable is introduced.

\begin{remark}[Collective continuation and separated convergence]
\label{loc:continuation-caution}
The original series is holomorphic on $\Omega$.  Since the finite
localized expression agrees with it on the nonempty subdomain
$\Omega_0$, the former gives a holomorphic continuation of the latter
collectively.  The degree estimate above does not prove that each
separated sector, or each individual exponent-shifted term, converges
on all of $\Omega$.  Finite cancellation between offset patterns may be
needed there.  This distinction is part of the theorem's domain
statement, not an appeal to formal rearrangement outside convergence.
\end{remark}

\subsection{Rank three: the complete two-index contact part}

For $r=3$, $d=4$, write $f_i(k)=(k+a_i)^{-s_i}$.  The unanchored
composition $(1,2)\mid(3)$ is described by
\[
 k_1\ge4,\qquad h=k_2-k_1\in\{-3,\ldots,4\},
 \qquad k_3-k_2>4.
\]
Summing the isolated third stencil first, exactly as in
Lemma~\ref{loc:singleton}, gives
\begin{align}
 C_3(k_1,k_2,k_3)&=24A(h),\label{loc:rank3-contact}\\
 A(h)&=\sum_{j,k=0}^4(-1)^{j+k}\binom4j\binom4k
                  (h+j-k)^2\ind_{\{h+j-k>0\}}.
                  \label{loc:rank3-A}
\end{align}
This entire integer contact table has a short generating expression:
\begin{equation}
 \sum_{h\in\Z}A(h)z^h
   =(1-z^{-1})^4(1-z)^4\frac{z(1+z)}{(1-z)^3}
   =z^{-3}(1+z)(1-z)^5.
 \label{loc:rank3-generating}
\end{equation}
Indeed, $z(1+z)/(1-z)^3=\sum_{h\ge1}h^2z^h$, and the two finite
stencils multiply this series by the displayed Laurent factors.
One may first work in $0<|z|<1$, where every intermediate series
converges.  The final Laurent polynomial proves
\[
\begin{array}{c|rrrrrrr}
 h&-3&-2&-1&0&1&2&3\\\hline
 A(h)&1&-4&5&0&-5&4&-1
\end{array}
\qquad A(h)=0\quad(|h|\ge4).
\]

The resulting two-index contribution of $(1,2)\mid(3)$ is
\begin{equation}
 \mathcal T_{12\mid3}
  =\frac1{24^2}\sum_{h=-3}^3 A(h)
       \sum_{0\le n<m}
       f_1(n+4)f_2(n+4+h)f_3(m+8+h).
 \label{loc:rank3-first}
\end{equation}
The other unanchored two-block composition, $(1)\mid(2,3)$, gives
\begin{equation}
 \mathcal T_{1\mid23}
  =\frac1{24^2}\sum_{h=-3}^3 A(h)
       \sum_{0\le n<m}
       f_1(n+4)f_2(m+8)f_3(m+8+h).
 \label{loc:rank3-second}
\end{equation}
Every remaining surviving sector has at most one free index.  Thus
\begin{equation}
 \mathcal Z_3=\mathcal T_{12\mid3}+\mathcal T_{1\mid23}
                         +\mathcal E_3,
 \label{loc:rank3-decomposition}
\end{equation}
where $\mathcal E_3$ is exactly the finite sum prescribed by
\eqref{loc:free-sector}--\eqref{loc:anchored-sector} over the remaining
compositions and anchor choices.  More explicitly, it consists of the
unanchored one-block sectors, all anchored one-block sectors, and the
anchored two-block sectors; anchored three-block sectors and
unanchored three-block sectors vanish by Lemma~\ref{loc:singleton}.
This is a finite formula for its one-index correction, with no
undetermined coefficient.

The contact contribution is not vacuous: the admissible tuple
$(5,6,20)$ has $h=1$ and
\begin{equation}
 C_3(5,6,20)=24A(1)=-120.
 \label{loc:rank3-nonzero}
\end{equation}
This proves the existence of two-index support in the construction.
Further identities can still simplify the complete sum for special
choices of the orders or shifts; \eqref{loc:rank3-nonzero} is not a
minimal-depth assertion.

\subsection{General integer stencils and stronger annihilation}

The support mechanism depends on polynomial moments of the stencil,
rather than on consecutive binomial coefficients.  The next result
retains $d=2r-2$, even when the stencil has a larger support.

\begin{corollary}[Finite integer stencils]
\label{loc:general-stencil}
Let $L\ge d$ be an integer and, for each label $i$, choose coefficients
$\alpha_{i,\lambda}\in\C$ supported on $0\le\lambda\le L$ such that
\begin{equation}
 \sum_{\lambda=0}^L\alpha_{i,\lambda}\lambda^k=0
                      \qquad(0\le k<d).
 \label{loc:stencil-moments}
\end{equation}
Replace \eqref{loc:weights} by
\[
 w_{i,\alpha}(n)=\sum_{\lambda=0}^L
                     \alpha_{i,\lambda}(n+a_i+\lambda)^{-s_i}.
\]
The corresponding ordered sum has the conclusions of
Theorem~\ref{loc:main}, with the same bounds
$\lfloor(r+1)/2\rfloor$ and $\lfloor r/2\rfloor$ on its free indices.
If every stencil also annihilates degree $d$, the bounds improve to
$\lfloor r/2\rfloor$ for the unanchored part and
$\lfloor(r-1)/2\rfloor$ for the anchored part.
\end{corollary}

\begin{proof}
The exact transported coefficient is now
\begin{equation}
 C_\alpha(\mathbf k)=
   \sum_{\boldsymbol\lambda\in\{0,\ldots,L\}^r}
          \left(\prod_i\alpha_{i,\lambda_i}\right)
          P(\mathbf k-\boldsymbol\lambda).
 \label{loc:general-coefficient}
\end{equation}
The factor $(d!)^{-r}$ is absent unless it was included in the chosen
coefficients.  Expanded absolute convergence on $\Omega_0$ justifies
the same transport.  A nonzero coefficient has adjacent gaps at least
$1-L$.  Make a cut at gaps greater than $L$, and anchor the first
block when $k_1<L$.

At a free singleton $i$, a degree-$d$ polynomial is reduced to its
leading coefficient times
\[
 \mu_i=\sum_{\lambda=0}^L
                 \alpha_{i,\lambda}(-\lambda)^d.
\]
This follows by expanding each power of $k_i-\lambda$ and using
\eqref{loc:stencil-moments}.  A second free singleton then occurs in
a polynomial of degree $d-2$, so its stencil annihilates the term.
The previous counting argument gives the original bounds.  If each
$\mu_i=0$, even the first free singleton annihilates the coefficient.
Every free block then uses at least two labels; with or without an
anchor this gives exactly the two improved bounds.

The full finite sector construction is unchanged after replacing the
geometric thresholds $d$ by $L$, the shift range by $\{0,\ldots,L\}$,
and the binomial stencil product by $\prod_i\alpha_{i,\lambda_i}$.
In particular internal offset increments belong to
$\{1-L,\ldots,L\}$, $c_\ell=L+u_{\max B_\ell}\ge b_\ell-1$, and
the unanchored first base is $L$.  The anchored first base has finitely
many choices $0\le h<L$, with the condition $\lambda_1\le h$.
Feasibility gives $u_i\ge i-\min B_\ell-L$, so the same base
substitutions keep all shifted arguments in the right half-plane.
The polynomial degree bound \eqref{loc:sector-degree} does not change.
Consequently $\Omega_0$, with $d=2r-2$, still suffices for convergence
of every separated sector, independently of $L$.

Finally Taylor's formula with a uniform remainder, applied to
$(n+a_i+\lambda)^{-s_i}$ through degree $d-1$ in $\lambda$, gives
\[
 w_{i,\alpha}(n)=O\bigl((1+n)^{-\Rea s_i-d}\bigr)
\]
locally uniformly in the parameters.  The polynomial terms vanish by
\eqref{loc:stencil-moments}; the finite support bounds the remainders.
Equation~\eqref{loc:vandermonde-bound} proves original normal
convergence on $\Omega$, and the collective continuation follows as
before.
\end{proof}

Every consecutive finite difference of order strictly greater than
$2r-2$ satisfies the stronger hypothesis.  For nonintegral rational
stencils, transport onto one integer lattice would instead require
separate residue classes; global moment cancellation alone does not
justify the integer-sector argument above.

The transport itself is discrete summation by parts.  Discrete residues
provide a classical framework for rational summability
\cite{chen-singer}.  Here the specific additional ingredient is the
support of the ordered chamber coefficient and the loss of two
Vandermonde degrees after a singleton is removed.  That mechanism works
with arbitrary complex orders and unequal shifts on the stated genuine
convergence domains, and resolves the finite-coordinate representation
question for this family in \cite{repo:finite}.
